% from ega4-2-fr.tex
\oldpage[IV-2]{97}

\begin{equation}
\dim(C_{\mathfrak r})+\operatorname{deg.tr.}_C B
\geq\dim(B_{\mathfrak q})+\operatorname{deg.tr.}_{k_1}k'.
\tag{5.5.8.2}\label{IV.2.s5.5.8.2-fr}
\end{equation}
D'autre part l'hypothèse de récurrence donne
\begin{equation}
\dim(A)+\operatorname{deg.tr.}_A C
\geq\dim(C_{\mathfrak r})+\operatorname{deg.tr.}_k k_1
\tag{5.5.8.3}\label{IV.2.s5.5.8.3-fr}
\end{equation}
d'où (5.5.8.1) en ajoutant (5.5.8.2) et (5.5.8.3) membre à membre. Pour
démontrer la seconde assertion, notons d'abord que toute sous-$A$-algèbre de
type fini de $C[T]$ est aussi une sous-$A$-algèbre de type fini de $B[T]$;
par récurrence sur $n$, on peut donc supposer que les deux membres de
(5.5.8.3) sont égaux. D'autre part, pour voir que les deux membres de
(5.5.8.2) sont égaux, il suffit, en vertu de (5.5.7), de vérifier que si
$\mathfrak n$ est un idéal maximal de $C_{\mathfrak r}[T]$ tel que
$\mathfrak n\cap C_{\mathfrak r}=\mathfrak rC_{\mathfrak r}$, alors
$(C_{\mathfrak r}[T])_{\mathfrak n}$ est caténaire; mais on a
$C_{\mathfrak r}[T]=S'^{-1}C[T]$ où $S'=C-\mathfrak r$; l'idéal
$\mathfrak n$ est donc de la forme $S'^{-1}\mathfrak n'$, où $\mathfrak n'$
est un idéal premier de $C[T]$ tel que $\mathfrak n'\cap C=\mathfrak r$,
d'où $(C_{\mathfrak r}[T])_{\mathfrak n}=(C[T])_{\mathfrak n'}$; si
$\mathfrak m'$ est un idéal maximal de $C[T]$ contenant $\mathfrak n'$,
$(C[T])_{\mathfrak n'}$ est donc un anneau local de l'anneau
$(C[T])_{\mathfrak m'}$, et comme par hypothèse ce dernier est caténaire, il
en est de même de $(C[T])_{\mathfrak n'}$ (0, 16.1.4).
\end{proof}

\subsection{Formule des dimensions et anneaux universellement caténaires}
\label{IV.2.s5.6-fr}

\begin{proposition}[5.6.1]
\EGAReaderTarget{ega4.c04.prop.5-6-1}{IV.2.s5.6.1-fr}
Soit $A$ un anneau noethérien. Les propriétés suivantes sont équivalentes:
\begin{enumerate}
\item[a)] Tout anneau de polynômes $A[T_1,\ldots,T_n]$ est caténaire.
\item[b)] Toute algèbre de type fini sur $A$ est caténaire.
\item[c)] $A$ est caténaire, et pour toute $A$-algèbre intègre locale,
essentiellement de type fini (1.3.8) $B'$, on a, en désignant par
$\mathfrak p$ l'image réciproque dans $A$ de l'idéal maximal $\mathfrak q$
de $B'$, par $A'$ l'image de $A_{\mathfrak p}$ dans $B'$, par $k$ et $k'$
les corps résiduels de $A'$ et $B'$ respectivement,
\begin{equation}
\dim(A')+\operatorname{deg.tr.}_{A'}(B')
=\dim(B')+\operatorname{deg.tr.}_k k'.
\tag{5.6.1.1}\label{IV.2.s5.6.1.1-fr}
\end{equation}
\end{enumerate}
\end{proposition}

Le fait que b) entraîne c) résulte de (5.5.8); en effet, $B'$ est une
$A'$-algèbre locale essentiellement de type fini (1.3.10), donc de la forme
$B''_{\mathfrak q''}$, où $B''$ est une sous-$A'$-algèbre de type fini de
$B'$, $\mathfrak q''$ un idéal premier de $B''$ au-dessus de l'idéal maximal
$\mathfrak p'$ de $A'$; en outre les corps des fractions de $B'$ et $B''$
sont les mêmes, donc
$\operatorname{deg.tr.}_{A'}(B')=\operatorname{deg.tr.}_{A'}(B'')$.
Pour démontrer (5.6.1.1), il suffit de montrer que (sous l'hypothèse b)) toute
sous-$A'$-algèbre $A_1$ de type fini de $B'[T]$ est caténaire; en effet les
deux membres de (5.5.8.1), où on remplace $A$, $B$ et $\mathfrak q$ par
$A'$, $B''$ et $\mathfrak q''$, seront alors égaux, d'où l'égalité (5.6.1.1).
Or l'hypothèse b) entraîne que toute $A$-algèbre essentiellement de type fini
est caténaire (0, 16.1.4), et $A_1$ est une telle $A$-algèbre (1.3.9).

Il est trivial que b) entraîne a); inversement, a) entraîne b), toute
$A$-algèbre de type fini étant quotient d'une algèbre de polynômes
(0, 16.1.4).

Reste à prouver que c) entraîne b). Comme tout quotient par un idéal premier
d'une $A$-algèbre de type fini est une $A$-algèbre intègre de type fini, il
s'agit de voir que, si $B$ est une $A$-algèbre intègre de type fini,
$\mathfrak q$ et $\mathfrak q'$ deux idéaux premiers de $B$ tels que
$\mathfrak q'\subset\mathfrak q$, on a (0, 16.1.4.2)
\begin{equation}
\dim(B_{\mathfrak q}/\mathfrak q'B_{\mathfrak q})
+\dim(B_{\mathfrak q'})=\dim(B_{\mathfrak q}).
\tag{5.6.1.2}\label{IV.2.s5.6.1.2-fr}
\end{equation}

